\begin{abstract}
To advance time series forecasting, a wide range of methods—from statistical techniques to deep learning architectures—have been proposed to improve prediction accuracy; however, most still follow a \textit{fast thinking} paradigm that maps historical patterns to future values with limited explicit reasoning or interpretable intermediate steps, undermining adaptability to evolving temporal dynamics.
Although slow-thinking LLMs exhibit strong multi-step reasoning, prompt-only solutions are constrained by cost, privacy, and insufficient domain alignment, motivating the training of LLMs to acquire \textit{slow thinking} time-series reasoning skills.
In this paper, we propose \Ours, a reinforcement fine-tuning framework designed to enhance the multi-step reasoning ability of LLMs for time series forecasting. 
Training starts with supervised fine-tuning for warmup adaptation, followed by reinforcement learning that improves the model’s generalization ability via reward-driven reasoning optimization.
Specifically, we introduce a fine-grained multi-objective reward to provide training signals tailored to time series forecasting, and GRIP (group-based relative importance for policy optimization) to leverage non-uniform sampling for encouraging the exploration of effective reasoning paths.
Experiments show that \Ours achieves competitive forecasting performance across diverse datasets
\footnote{Code is at \url{https://github.com/ustc-time-series/Time-R1}}.
\end{abstract}
% To advance time series forecasting (TSF), various methods have been proposed to improve prediction accuracy, ranging from statistical techniques to data-driven deep learning architectures.
% Despite their effectiveness, most existing methods still adhere to a \textit{fast thinking} paradigm-relying on extracting historical patterns and mapping them to future values as their core modeling philosophy, lacking an explicit reasoning process. Meanwhile, emerging slow-thinking LLMs (e.g., OpenAI-o1) have shown remarkable multi-step reasoning capabilities, offering an alternative way to overcome these issues. However, prompt engineering alone presents several limitations—including high computational cost, privacy risks, and limited capacity for in-depth domain-specific time series reasoning. To address these limitations, a more promising approach is to train LLMs to develop \textit{slow thinking} capabilities and acquire strong time series reasoning skills.
% For this purpose, we propose Time-R1, a two-stage reinforcement fine-tuning framework designed to enhance multi-step reasoning ability of LLMs for time series forecasting. Specifically, the first stage conducts supervised fine-tuning for warmup adaptation, while the second stage employs reinforcement learning to improve the model's generalization ability. Particularly, we design a fine-grained multi-objective reward specifically for time series forecasting, and then introduce GRIP (group-based relative importance for policy optimization), which leverages non-uniform sampling to further encourage and optimize the model's exploration of effective reasoning paths.
% Experiments shows Time-R1 significantly improves forecast performance in diverse datasets.
% \footnote{The code is at  \url{https://anonymous.4open.science/r/Time-R1-ICML-2026}.}.

% In this paper, we propose Time-R1, a reinforcement fine-tuning framework designed to enhance the multi-step reasoning ability of LLMs for time series forecasting through a fine-grained multi-objective reward and GRIP (group-based relative importance for policy optimization).